% FujiCV manuscript skeleton.
% Compiles with pdflatex out of the box (article class).
% For SIVP (Springer), switch the class to: \documentclass[smallextended]{svjour3}
% and add \journalname{Signal, Image and Video Processing}. For SoftwareX, use
% the Elsevier elsarticle class. Section content is portable across all three.
\documentclass[11pt]{article}

\usepackage[utf8]{inputenc}
\usepackage[margin=1in]{geometry}
\usepackage{graphicx}
\usepackage{booktabs}
\usepackage{amsmath}
\usepackage{listings}
\usepackage{xcolor}
\usepackage{hyperref}
\usepackage{url}

\lstset{basicstyle=\ttfamily\small, breaklines=true, frame=single,
        backgroundcolor=\color{gray!8}, columns=fullflexible}

\title{FujiCV: A Lightweight, Batteries-Included Library for Image
Classification, Regression, and Retrieval in PyTorch}

\author{Sabarinathan\\
\small Independent Researcher \\
\small \texttt{sabarinathantce@gmail.com}}

\date{\today}

\begin{document}
\maketitle

\begin{abstract}
Fine-tuning image models in PyTorch repeatedly requires the same scaffolding:
data loading, augmentation, a mixed-precision training loop, checkpoint
selection, metric tracking, distributed handling, and an export path. Writing
this by hand is verbose and error-prone, while heavier frameworks trade
boilerplate for their own abstractions and are oriented mainly toward
classification. We present \textbf{FujiCV}, an open-source library that provides
a single, config-light \texttt{ModelBuilder}\,$\rightarrow$\,\texttt{Trainer}
pipeline spanning classification, regression, multi-label, and
metric-learning/retrieval on top of the \texttt{timm}, \texttt{torchvision}, and
Hugging~Face model zoos. FujiCV includes modern training techniques (EMA, SWA,
SAM, Model Soups, Mixup/CutMix, LLRD, gradient accumulation, DDP,
\texttt{torch.compile}), explainability (Grad-CAM), hyperparameter optimization,
retrieval (ArcFace/CosFace/GeM), and deployment (ONNX, TorchScript, post-training
quantization) out of the box. We show that FujiCV substantially reduces the code
required for a complete training pipeline while adding negligible throughput and
memory overhead relative to an equivalent hand-written PyTorch loop, and that the
abstraction preserves accuracy across standard benchmarks. FujiCV is released
under the Apache-2.0 license and is available on PyPI and GitHub.
\end{abstract}

\textbf{Keywords:} deep learning; computer vision; image classification; transfer
learning; PyTorch; open-source software

\section{Introduction}
\label{sec:intro}
% - The recurring boilerplate problem when fine-tuning image models.
% - Correctness pitfalls (distributed metric aggregation, checkpoint races,
%   normalization mismatch) that silently degrade results.
% - The gap between raw PyTorch (control, boilerplate) and heavy frameworks
%   (abstraction, learning curve, classification-centric).
% - Contributions (enumerated below).

\paragraph{Contributions.}
\begin{enumerate}
  \item A unified, config-light API covering \emph{classification, regression,
  multi-label, and retrieval} in one interface on top of \texttt{timm},
  \texttt{torchvision}, and Hugging~Face backbones.
  \item Modern training, explainability, retrieval, and deployment techniques
  available out of the box with minimal code.
  \item An empirical study showing substantial boilerplate reduction at
  negligible compute overhead and no accuracy loss.
  \item An extensively tested, continuously integrated, reproducible
  open-source release.
\end{enumerate}

\section{Related Work}
\label{sec:related}
% Position vs. raw PyTorch, PyTorch Lightning, fastai, Keras, timm (backbones
% only), MMPretrain, Ludwig. Emphasize: smaller surface area, broader task
% coverage, integrated deployment. See Table~\ref{tab:comparison}.

\begin{table}[t]
\centering
\caption{Positioning of FujiCV relative to related frameworks. ``Tasks'' denotes
first-class support (C: classification, R: regression, ML: multi-label, Ret:
retrieval).}
\label{tab:comparison}
\small
\begin{tabular}{lccccc}
\toprule
Framework & Tasks & Backbones & Built-in advanced training & Deployment & Retrieval \\
\midrule
Raw PyTorch          & ---        & any            & manual         & manual        & manual \\
PyTorch Lightning    & C          & any            & partial        & partial       & no     \\
fastai               & C, R       & timm/tv        & yes            & partial       & no     \\
Keras                & C, R       & Keras zoo      & partial        & yes           & no     \\
\textbf{FujiCV}      & C, R, ML, Ret & timm/tv/HF  & yes            & yes (ONNX/TS/PTQ) & yes \\
\bottomrule
\end{tabular}
\end{table}

\section{Design and Architecture}
\label{sec:design}
% Figure: pipeline diagram ModelBuilder -> Trainer -> eval/export.
% Philosophy: plain Python constructors, no hidden global config; every knob is
% a keyword argument; resolved config (with package+git version) saved per run.

\begin{lstlisting}[language=Python, caption={A complete CIFAR-10 pipeline in FujiCV.}, label={lst:quickstart}]
from torch.utils.data import DataLoader
import torch
from fujicv.data.datasets import get_default_dataset
from fujicv.data.transforms import get_train_transforms, get_val_transforms
from fujicv.engine.trainer import Trainer
from fujicv.losses.classification import CrossEntropyLoss
from fujicv.metrics.classification import Accuracy
from fujicv.models.builder import ModelBuilder

train_ds, val_ds, c2i = get_default_dataset(
    "cifar10", root="data",
    train_transform=get_train_transforms(32), val_transform=get_val_transforms(32))
train_loader = DataLoader(train_ds, batch_size=128, shuffle=True,  num_workers=2)
val_loader   = DataLoader(val_ds,   batch_size=128, shuffle=False, num_workers=2)

model = ModelBuilder("resnet18", task="classification", num_outputs=10,
                     image_size=32).build()
trainer = Trainer(model=model, train_loader=train_loader, val_loader=val_loader,
                  loss_fn=CrossEntropyLoss(), metrics={"accuracy": Accuracy()},
                  optimizer=torch.optim.AdamW(model.parameters(), lr=1e-3),
                  epochs=10, task="classification", output_dir="runs/cifar10",
                  monitor_metric="val_accuracy")
history = trainer.train()   # -> best.pt  last.pt  history.csv  history.json
\end{lstlisting}

\section{Features}
\label{sec:features}
% Narrative over Table~\ref{tab:features}. Highlight the correctness work:
% rank-guarded DDP checkpoints, all-gathered DDP metrics, crash-resilient history.

\begin{table}[t]
\centering
\caption{Feature coverage in FujiCV.}
\label{tab:features}
\small
\begin{tabular}{ll}
\toprule
Category & Included \\
\midrule
Backbones        & timm, torchvision, Hugging~Face (incl. CLIP/SigLIP/DINOv2 towers) \\
Tasks / heads    & classification, regression, multi-label, ArcFace/CosFace/Sub-center \\
Pooling          & average, max, GeM, attention \\
Losses           & 15+ (Focal, LabelSmoothing, CORAL/CORN, Quantile, Asymmetric, \ldots) \\
Metrics          & 16+ (Accuracy, F1, AUROC, mAP, MAE, RMSE, R$^2$, \ldots) \\
Augmentation     & Albumentations presets, RandAugment, Mixup, CutMix \\
Training         & AMP, grad clip, grad accumulation, EMA, SWA, SAM, early stopping \\
LR utilities     & LR finder, cosine warmup, OneCycle, LLRD, layer freezing \\
Multi-GPU        & DistributedDataParallel (rank-guarded ckpt, all-gathered metrics) \\
Weight averaging & Model soups (uniform, greedy) \\
Explainability   & Grad-CAM, Grad-CAM++, confusion matrix, calibration \\
Retrieval        & embedding extraction, cosine/FAISS kNN, Recall@K, mAP@K \\
HPO              & Optuna search with pruning \\
Export           & ONNX (+INT8), TorchScript, dynamic/static PTQ \\
Logging          & W\&B, TensorBoard, MLflow (pluggable \texttt{BaseLogger}) \\
\bottomrule
\end{tabular}
\end{table}

\section{Experiments}
\label{sec:experiments}
Hardware, software versions, and seeds for all experiments are reported in the
supplementary material and captured automatically by the benchmark harness
(\texttt{paper/benchmarks/}). Every table below is regenerated from JSON results
by \texttt{make\_tables.py}.

\subsection{E1: Boilerplate reduction (lines of code)}
\label{sec:e1}
For an identical pipeline (AMP, early stopping, checkpointing, metric tracking)
on CelebA attribute classification, we count non-comment, non-blank lines of the
training script (Table~\ref{tab:loc}).

\begin{table}[h]
\centering
\caption{Lines of code for an equivalent training pipeline (lower is better).}
\label{tab:loc}
\small
\begin{tabular}{lcc}
\toprule
Framework & LoC & Relative \\
\midrule
Raw PyTorch       & 152 & 1.00$\times$ \\
PyTorch Lightning & \textit{TBD} & \textit{TBD} \\
fastai            & \textit{TBD} & \textit{TBD} \\
\textbf{FujiCV}   & 89  & \textbf{0.59$\times$} \\
\bottomrule
\end{tabular}
\end{table}

\subsection{E2: Throughput and memory overhead}
\label{sec:e2}
We measure training throughput (images/s) and peak GPU memory against an
equivalent hand-written PyTorch loop on identical data and model, reporting the
mean $\pm$ standard deviation over $N{=}3$ runs after a warmup epoch, on a single
GPU and with 2-GPU DistributedDataParallel (Table~\ref{tab:throughput}).

\begin{table}[h]
\centering
\caption{Throughput and peak memory (ResNet-18, CelebA, $64{\times}64$). Values
are placeholders pending the final controlled runs; see \S\ref{sec:threats}.}
\label{tab:throughput}
\small
\begin{tabular}{llccc}
\toprule
Setting & Framework & Throughput (img/s) $\uparrow$ & Peak mem (MB) $\downarrow$ & Overhead \\
\midrule
Single-GPU & Raw PyTorch & \textit{TBD} & \textit{TBD} & --- \\
Single-GPU & FujiCV      & \textit{TBD} & \textit{TBD} & \textit{TBD} \\
2$\times$ DDP & Raw PyTorch & \textit{TBD} & \textit{TBD} & --- \\
2$\times$ DDP & FujiCV      & \textit{TBD} & \textit{TBD} & \textit{TBD} \\
\bottomrule
\end{tabular}
\end{table}

\subsection{E3: Accuracy parity}
\label{sec:e3}
To confirm the abstraction does not cost accuracy, we reproduce standard results
on CIFAR-10/100 and Oxford-IIIT Pet and compare against a hand-written loop and
published baselines (Table~\ref{tab:accuracy}).

\begin{table}[h]
\centering
\caption{Top-1 validation accuracy (\%). FujiCV vs. an equivalent hand-written
loop.}
\label{tab:accuracy}
\small
\begin{tabular}{lccc}
\toprule
Dataset (model) & Raw PyTorch & FujiCV & $\Delta$ \\
\midrule
CIFAR-10 (ResNet-18)  & \textit{TBD} & \textit{TBD} & \textit{TBD} \\
CIFAR-100 (ResNet-50) & \textit{TBD} & \textit{TBD} & \textit{TBD} \\
Oxford Pet (ResNet-50)& \textit{TBD} & \textit{TBD} & \textit{TBD} \\
\bottomrule
\end{tabular}
\end{table}

\subsection{E4: Value of built-in techniques (ablation)}
\label{sec:e4}
We enable each built-in technique in turn to quantify the accuracy it buys
``for free'' from the user's perspective (Table~\ref{tab:ablation}).

\begin{table}[h]
\centering
\caption{Effect of built-in techniques on validation accuracy (single dataset,
fixed budget). Deltas are relative to the baseline.}
\label{tab:ablation}
\small
\begin{tabular}{lcc}
\toprule
Configuration & Val acc (\%) & $\Delta$ \\
\midrule
Baseline           & \textit{TBD} & --- \\
\ + Mixup/CutMix   & \textit{TBD} & \textit{TBD} \\
\ + EMA            & \textit{TBD} & \textit{TBD} \\
\ + SWA            & \textit{TBD} & \textit{TBD} \\
\ + Model soup     & \textit{TBD} & \textit{TBD} \\
\bottomrule
\end{tabular}
\end{table}

\subsection{E5: Deployment (quantization)}
\label{sec:e5}
We report model size and CPU inference latency for FP32, dynamic INT8, and
static (FX) INT8 post-training quantization (Table~\ref{tab:deploy}).

\begin{table}[h]
\centering
\caption{Post-training quantization: size and CPU latency.}
\label{tab:deploy}
\small
\begin{tabular}{lccc}
\toprule
Variant & Size (MB) $\downarrow$ & Latency (ms) $\downarrow$ & Acc.\ drop \\
\midrule
FP32          & \textit{TBD} & \textit{TBD} & --- \\
Dynamic INT8  & \textit{TBD} & \textit{TBD} & \textit{TBD} \\
Static INT8   & \textit{TBD} & \textit{TBD} & \textit{TBD} \\
\bottomrule
\end{tabular}
\end{table}

\subsection{E6: Retrieval}
\label{sec:e6}
Using an ArcFace head with GeM pooling, we report Recall@$K$ and mAP@$K$ on a
standard retrieval benchmark to demonstrate the breadth of the API
(Table~\ref{tab:retrieval}).

\begin{table}[h]
\centering
\caption{Retrieval performance (ArcFace + GeM).}
\label{tab:retrieval}
\small
\begin{tabular}{lccc}
\toprule
Dataset & Recall@1 & Recall@5 & mAP@10 \\
\midrule
CUB-200 & \textit{TBD} & \textit{TBD} & \textit{TBD} \\
\bottomrule
\end{tabular}
\end{table}

\section{Software Quality and Reproducibility}
\label{sec:quality}
% Test suite size and coverage; CI matrix (Python 3.9/3.10/3.11); strict docs
% build; resolved-config capture (package + git hash); Zenodo DOI for the tagged
% release; PyPI distribution.

\section{Threats to Validity}
\label{sec:threats}
% Be explicit and honest: (i) throughput comparisons are sensitive to data
% loading, worker count, and shared-machine contention -- report warmup, fixed
% workers, and multiple runs; (ii) many components wrap timm/albumentations/HF,
% so the contribution is standardized, tested integration rather than new
% algorithms; (iii) LoC is an imperfect proxy for effort, measured with a fixed
% counting rule.

\section{Limitations and Conclusion}
\label{sec:conclusion}
% Summary; current limitations (single-node focus for DDP validation, image
% tasks only); future work.

\section*{Availability}
Source code: \url{https://github.com/dsabarinathan/fujicv}. Package:
\url{https://pypi.org/project/fujicv/}. License: Apache-2.0.

\bibliographystyle{plain}
\bibliography{paper}

\end{document}
